\documentclass[11pt,a4paper]{article}

% ── Packages ──────────────────────────────────────────────────────────────────
\usepackage[utf8]{inputenc}
\usepackage[T1]{fontenc}
\usepackage[french]{babel}
\usepackage{amsmath,amssymb,amsthm,mathtools}
\usepackage{bm}
\usepackage{geometry}
\usepackage{hyperref}
\usepackage{xcolor}
\usepackage{algorithm}
\usepackage{algpseudocode}
\usepackage{booktabs}
\usepackage{cleveref}
\usepackage{thmtools}

\geometry{margin=2.5cm}
\hypersetup{colorlinks=true, linkcolor=blue!70!black, citecolor=green!50!black}

% ── Theorem environments ───────────────────────────────────────────────────────
\theoremstyle{plain}
\newtheorem{theorem}{Théorème}[section]
\newtheorem{lemma}[theorem]{Lemme}
\newtheorem{corollary}[theorem]{Corollaire}
\newtheorem{proposition}[theorem]{Proposition}

\theoremstyle{definition}
\newtheorem{definition}[theorem]{Définition}
\newtheorem{assumption}[theorem]{Hypothèse}
\newtheorem{remark}[theorem]{Remarque}
\newtheorem{example}[theorem]{Exemple}

\theoremstyle{remark}
\newtheorem*{notation}{Notation}

% ── Custom macros ──────────────────────────────────────────────────────────────
\newcommand{\R}{\mathbb{R}}
\newcommand{\N}{\mathbb{N}}
\newcommand{\E}{\mathbb{E}}
\newcommand{\Prob}{\mathbb{P}}
\newcommand{\calX}{\mathcal{X}}
\newcommand{\calT}{\mathcal{T}}
\newcommand{\calF}{\mathcal{F}}
\newcommand{\calB}{\mathcal{B}}
\newcommand{\calL}{\mathcal{L}}
\newcommand{\logZ}{\log Z}
\newcommand{\pF}{P_F}
\newcommand{\pB}{P_B}
\newcommand{\pth}{P_\theta}
\newcommand{\pR}{p_R}
\newcommand{\lq}{\ell_q}
\newcommand{\TV}{\mathrm{TV}}
\newcommand{\argmin}{\operatorname{arg\,min}}
\newcommand{\supp}{\operatorname{supp}}
\newcommand{\norm}[1]{\left\lVert#1\right\rVert}

% ── Coloured boxes for key results ────────────────────────────────────────────
\usepackage{mdframed}
\newmdenv[linecolor=blue!40,backgroundcolor=blue!5,
          linewidth=1.5pt,roundcorner=4pt]{mainbox}
\newmdenv[linecolor=green!50!black,backgroundcolor=green!5,
          linewidth=1.5pt,roundcorner=4pt]{corrbox}

% ══════════════════════════════════════════════════════════════════════════════
\begin{document}

\title{%
  \textbf{Preuve formelle : GFlowNet avec mesure de base structurée}\\[4pt]
  \large Base-Measure Trajectory Balance — Théorèmes, Lemmes et Protocole de validation
}
\author{%
  \textit{Document de travail}
}
\date{\today}
\maketitle
\tableofcontents
\vspace{1em}

% ══════════════════════════════════════════════════════════════════════════════
\section{Cadre et notations}
\label{sec:setup}

% ── 1.1 Espace d'états et trajectoires ───────────────────────────────────────
\subsection{Espace d'états et trajectoires}

Soit $\calX$ un ensemble fini d'états terminaux, $|\calX| = N < \infty$.
Un \emph{GFlowNet} génère des états $x \in \calX$ séquentiellement en $D$ étapes :

\begin{equation}
  \tau = (s_0, a_1, s_1, a_2, \ldots, a_D, x),
  \qquad x = s_D \in \calX,
\end{equation}

où $s_0$ est l'état initial vide, $a_t \in \{0,\ldots,K-1\}$ est l'action à l'étape $t$,
et la transition $s_t = f(s_{t-1}, a_t)$ est déterministe.
On note $\calT(x)$ l'ensemble des trajectoires menant à $x$.

\begin{assumption}[Injectivité des trajectoires]
\label{asm:inject}
La fonction $(a_1,\ldots,a_D) \mapsto x$ est injective : chaque état terminal
possède exactement une trajectoire, i.e., $|\calT(x)| = 1$ pour tout $x$.
\end{assumption}

\begin{remark}
L'hypothèse~\ref{asm:inject} couvre les espaces séquentiels sans cycle utilisés
dans les expériences : $\{0,\ldots,K-1\}^D$ avec encodage position-valeur.
Elle simplifie la marginalistion $\pF(\tau) = \pth(x)$ (voir Lemme~\ref{lem:marg}).
La preuve du Théorème~\ref{thm:main} se généralise aux DAG multi-trajectoires
par le même argument variationnel, en remplaçant la somme sur $\calT(x)$ par
la somme sur tous les chemins racine$\to x$.
\end{remark}

% ── 1.2 Politique forward et distribution marginale ──────────────────────────
\subsection{Politique forward et distribution marginale}

Soit $\theta \in \R^p$ les paramètres d'un réseau de neurones.
La \emph{politique forward} est une famille de catégorielles :
\begin{equation}
  \pF^\theta(a \mid s) = \operatorname{softmax}\bigl(f_\theta(s)\bigr)_a
  \;\in\; (0,1), \qquad \forall s, \forall a.
\end{equation}

La probabilité d'une trajectoire complète est :
\begin{equation}
  \pF^\theta(\tau) = \prod_{t=1}^{D} \pF^\theta(a_t \mid s_{t-1}).
\end{equation}

Sous l'Hypothèse~\ref{asm:inject}, la \emph{distribution marginale sur $\calX$} est :
\begin{equation}
  \pth(x) = \pF^\theta(\tau_x), \qquad \text{où } \tau_x \text{ est l'unique trajectoire menant à } x.
\end{equation}

\begin{assumption}[Positivité]
\label{asm:pos}
Pour tout $\theta$, $\pF^\theta(a \mid s) > 0$ pour tout couple $(s,a)$.
Ceci garantit $\pth(x) > 0$ pour tout $x \in \calX$.
\end{assumption}

% ── 1.3 Récompense et prior ───────────────────────────────────────────────────
\subsection{Récompense, mesure de base et distributions cibles}

\begin{definition}[Récompense et distribution cible]
Soit $R : \calX \to \R_{>0}$ une fonction de récompense strictement positive.
La distribution cible \emph{standard} est :
\begin{equation}
  \pR(x) = \frac{R(x)}{Z_R}, \qquad Z_R = \sum_{x' \in \calX} R(x').
\end{equation}
\end{definition}

\begin{definition}[Mesure de base structurée]
\label{def:base}
Soit $q : \calX \to \R_{>0}$ une \emph{mesure de base} (prior structuré),
normalisée : $\sum_{x} q(x) = 1$.
La distribution cible \emph{modifiée} est :
\begin{equation}
  p_{q,R}(x) = \frac{q(x)\,R(x)}{Z_q}, \qquad Z_q = \sum_{x'} q(x')\,R(x').
\end{equation}
\end{definition}

\begin{definition}[Politique comportementale]
\label{def:bpol}
Soit $\beta : \calX \to \R_{>0}$ une \emph{politique comportementale} (behavioral policy),
$\sum_x \beta(x) = 1$, indépendante de $\theta$. Elle guide l'exploration (replay,
sampling biaisé) mais n'entre pas dans la loss.
\end{definition}

% ── 1.4 Partition log-Z ──────────────────────────────────────────────────────
\subsection{Partition logarithmique}

Le scalaire $\logZ \in \R$ est un paramètre appris séparément de $\theta$
(typiquement, un \texttt{nn.Parameter} initialisé à 0).

% ══════════════════════════════════════════════════════════════════════════════
\section{Objectif Trajectory Balance avec mesure de base}
\label{sec:objective}

\subsection{Loss TB standard}

La loss \emph{Trajectory Balance} (Malkin et al., 2022) est :
\begin{equation}
  \calL_{\mathrm{TB}}(\theta, Z)
  = \E_{\tau \sim \beta}\Bigl[
      \bigl(\logZ + \log\pF^\theta(\tau) - \log R(x_\tau)\bigr)^2
    \Bigr],
\end{equation}
où $x_\tau = s_D(\tau)$ est l'état terminal de $\tau$.
Sous l'Hypothèse~\ref{asm:inject} :
$\log\pF^\theta(\tau) = \log\pth(x_\tau)$.

\subsection{Loss TB avec mesure de base (TB+q)}

\begin{definition}[Loss TB+q]
\label{def:tbq}
La loss \emph{TB avec mesure de base} est :
\begin{equation}
  \boxed{
  \calL_{\mathrm{TB+q}}(\theta, Z)
  = \E_{\tau \sim \beta}\Bigl[
      \bigl(\logZ + \log\pF^\theta(\tau) - \log q(x_\tau) - \log R(x_\tau)\bigr)^2
    \Bigr].
  }
\end{equation}
\end{definition}

\begin{remark}[Interprétation]
Par rapport à TB standard, on substitue $\log R(x)$ par $\log q(x) + \log R(x)$,
i.e., on impose au GFlowNet d'apprendre la densité $q(x)\,R(x)$ plutôt que $R(x)$.
\end{remark}

% ══════════════════════════════════════════════════════════════════════════════
\section{Lemmes auxiliaires}
\label{sec:lemmas}

% ── Lemme 1 : Marginalisaton ──────────────────────────────────────────────────
\begin{lemma}[Marginalisaton des trajectoires]
\label{lem:marg}
Sous l'Hypothèse~\ref{asm:inject}, pour tout $\theta$ :
\begin{equation}
  \sum_{\tau : x_\tau = x} \pF^\theta(\tau) = \pth(x) \qquad \forall x \in \calX.
\end{equation}
\end{lemma}

\begin{proof}
Sous l'injectivité, $|\calT(x)| = 1$, donc la somme est réduite à un seul terme :
$\sum_{\tau \in \calT(x)} \pF^\theta(\tau) = \pF^\theta(\tau_x) = \pth(x)$.
\end{proof}

% ── Lemme 2 : Zéro de la loss par état ───────────────────────────────────────
\begin{lemma}[Condition d'annulation terme à terme]
\label{lem:zero}
Soit $(\theta^*, Z^*)$ tel que $\calL_{\mathrm{TB+q}}(\theta^*, Z^*) = 0$.
Alors, pour $\beta$-presque tout $x \in \calX$ :
\begin{equation}
  \log Z^* + \log \pth[*](x) = \log q(x) + \log R(x).
\end{equation}
\end{lemma}

\begin{proof}
La loss est une espérance de carrés, donc nulle si et seulement si chaque terme
est nul $\beta$-p.s. Sous l'Hypothèse~\ref{asm:pos}, $\supp(\beta) = \calX$,
donc la condition vaut pour tout $x$.
\end{proof}

% ── Lemme 3 : Détermination de Z* ────────────────────────────────────────────
\begin{lemma}[Valeur de la partition optimale]
\label{lem:Zstar}
Si $\calL_{\mathrm{TB+q}}(\theta^*, Z^*) = 0$, alors :
\begin{equation}
  Z^* = Z_q = \sum_{x \in \calX} q(x)\,R(x).
\end{equation}
\end{lemma}

\begin{proof}
Du Lemme~\ref{lem:zero}, $\pth[*](x) = Z^{*-1} q(x) R(x)$ pour tout $x$.
En sommant sur $\calX$ et en utilisant $\sum_x \pth[*](x) = 1$ :
\begin{equation}
  1 = \frac{1}{Z^*} \sum_{x} q(x) R(x) = \frac{Z_q}{Z^*},
\end{equation}
d'où $Z^* = Z_q$.
\end{proof}

% ── Lemme 4 : Unicité du point fixe ──────────────────────────────────────────
\begin{lemma}[Unicité du point fixe $\pth$]
\label{lem:unique}
L'annulation de $\calL_{\mathrm{TB+q}}$ contraint $\pth$ de manière unique :
il existe une unique distribution $p^* \in \Delta(\calX)$ telle que
$\calL_{\mathrm{TB+q}}(\theta, Z) = 0$ implique $\pth = p^*$.
\end{lemma}

\begin{proof}
Des Lemmes~\ref{lem:zero} et \ref{lem:Zstar}, tout optimum global vérifie :
\begin{equation}
  \pth(x) = \frac{q(x)\,R(x)}{Z_q} \qquad \forall x \in \calX.
\end{equation}
Cette expression est indépendante de $\theta$ — elle ne dépend que de $q$ et $R$.
Donc $p^* = p_{q,R}$ est l'unique point fixe.
\end{proof}

% ── Lemme 5 : L'espérance sous β est non biaisée ─────────────────────────────
\begin{lemma}[Invariance du point fixe par rapport à $\beta$]
\label{lem:beta}
Pour toute politique comportementale $\beta$ de support plein ($\beta(x) > 0$
pour tout $x$), le point fixe de $\calL_{\mathrm{TB+q}}$ est identique :
$\pth = p_{q,R}$.
\end{lemma}

\begin{proof}
La condition d'annulation (Lemme~\ref{lem:zero}) doit tenir pour tout $x$
dans le support de $\beta$. Si $\supp(\beta) = \calX$ (complet), elle tient
pour tout $x \in \calX$ sans dépendre de la valeur de $\beta(x)$.
Le minimum est atteint au même $p^*$ indépendamment de $\beta$.

Formellement, soit $h(x) = \logZ + \log\pth(x) - \log q(x) - \log R(x)$.
La loss vaut $\sum_x \beta(x) h(x)^2$.
Cette somme pondérée est nulle ssi $h(x) = 0$ pour tout $x \in \supp(\beta)$.
Avec support plein, $h \equiv 0$, et la solution ne dépend pas des poids $\beta(x)$.
\end{proof}

% ══════════════════════════════════════════════════════════════════════════════
\section{Théorèmes principaux}
\label{sec:theorems}

% ── Théorème 1 ────────────────────────────────────────────────────────────────
\begin{mainbox}
\begin{theorem}[TB avec mesure de base — Point fixe]
\label{thm:main}
Soient $R : \calX \to \R_{>0}$ et $q : \calX \to \R_{>0}$, $\sum_x q(x) = 1$.
Sous les Hypothèses~\ref{asm:inject} et~\ref{asm:pos},
tout minimiseur global $(\theta^*, Z^*)$ de $\calL_{\mathrm{TB+q}}$ satisfait :
\begin{equation}
  \pth[*](x) = \frac{q(x)\,R(x)}{Z_q},
  \qquad Z_q = \sum_{x' \in \calX} q(x')\,R(x'),
  \qquad \forall x \in \calX.
\end{equation}
En particulier, le minimum global vaut zéro : $\min \calL_{\mathrm{TB+q}} = 0$.
\end{theorem}
\end{mainbox}

\begin{proof}
\textbf{(Existence d'un minimiseur à valeur zéro.)}
Posons $p^*(x) = q(x)R(x)/Z_q$ et $Z^* = Z_q$.
Par construction $\sum_x p^*(x) = 1$.
L'Hypothèse~\ref{asm:pos} assure que $p^*$ est dans l'intérieur du simplexe,
donc réalisable par un réseau softmax (universal approximation pour les réseaux
de largeur suffisante, ou directement en paramétrant $\log\pth(x)$ par un
vecteur de $\R^N$). Pour ce choix :
\begin{equation}
  \logZ^* + \log\pth[*](x) - \log q(x) - \log R(x)
  = \log Z_q + \log\frac{q(x)R(x)}{Z_q} - \log q(x) - \log R(x) = 0,
\end{equation}
donc $\calL_{\mathrm{TB+q}}(\theta^*, Z^*) = 0$.

\textbf{(Unicité du point fixe.)}
Réciproquement, si $\calL_{\mathrm{TB+q}}(\theta, Z) = 0$,
le Lemme~\ref{lem:zero} donne $\log\pth(x) = \log q(x) + \log R(x) - \logZ$
pour tout $x$, et le Lemme~\ref{lem:Zstar} impose $Z = Z_q$.
Donc $\pth = p^*$ est uniquement déterminé.
\end{proof}

% ── Corollaire 1 ──────────────────────────────────────────────────────────────
\begin{corrbox}
\begin{corollary}[Neutralisation asymptotique du prior comportemental]
\label{cor:behav}
Sous les hypothèses du Théorème~\ref{thm:main}, si l'on utilise la loss
TB standard $\calL_{\mathrm{TB}}$ (sans $\log q$) avec une politique
comportementale $\beta$ quelconque de support plein, alors :
\begin{equation}
  \pth[*](x) = \frac{R(x)}{Z_R}, \qquad \forall x \in \calX,
\end{equation}
indépendamment de $\beta$.
\end{corollary}
\end{corrbox}

\begin{proof}
Appliquons le Théorème~\ref{thm:main} avec $q \equiv 1/N$ (prior uniforme).
Alors $q(x)\,R(x)/Z_q = R(x)/(N \cdot Z_q)$ avec $Z_q = \sum_x R(x)/N = Z_R/N$.
On obtient $\pth[*](x) = R(x)/Z_R = \pR(x)$, indépendamment de $q = 1/N$.

Plus directement : la loss TB standard s'écrit $\calL_{\mathrm{TB}}(\theta,Z)
= \E_\beta[(\logZ + \log\pth(x) - \log R(x))^2]$.
Par le Lemme~\ref{lem:beta} (avec $q \equiv 1$, ou en absorbant $q$ dans $Z$),
le point fixe est $\pth(x) \propto R(x)$, sans dépendre de $\beta$.

\textbf{Interprétation.}
Le prior comportemental $\beta$ modifie la \emph{vitesse de convergence}
(couverture, variance d'estimation) mais pas le \emph{point fixe} asymptotique.
C'est la distinction fondamentale entre Niveau~1 (mesure de base dans la loss)
et Niveau~2 (prior dans l'exploration seulement).
\end{proof}

% ── Corollaire 2 ──────────────────────────────────────────────────────────────
\begin{corrbox}
\begin{corollary}[Sampling uniforme avec mesure de base]
\label{cor:uniform}
Sous les hypothèses du Théorème~\ref{thm:main}, si l'on utilise la loss TB+q
avec un sampling uniforme $\beta = \mathrm{Unif}(\calX)$, alors :
\begin{equation}
  \pth[*](x) = \frac{q(x)\,R(x)}{Z_q}, \qquad \forall x \in \calX.
\end{equation}
\end{corollary}
\end{corrbox}

\begin{proof}
Cas particulier du Théorème~\ref{thm:main} et du Lemme~\ref{lem:beta}
avec $\beta = \mathrm{Unif}(\calX)$, qui a support plein.
\end{proof}

\begin{remark}[Portée pratique]
Le Corollaire~\ref{cor:uniform} montre que même un sampling naïf (uniforme)
suffit pour que la mesure de base $q$ soit incorporée correctement dans la
distribution apprise — à condition qu'elle soit dans la loss, pas seulement
dans la politique comportementale.
\end{remark}

% ── Proposition sur la borne de TV ────────────────────────────────────────────
\begin{proposition}[Borne TV à l'optimum approché]
\label{prop:tv}
Soit $\varepsilon > 0$ tel que $\calL_{\mathrm{TB+q}}(\theta, Z) \leq \varepsilon^2$.
Alors, sous l'Hypothèse~\ref{asm:pos} :
\begin{equation}
  \TV(\pth, p_{q,R})
  \;\leq\;
  \frac{\sqrt{N\,\varepsilon}}{2\,\min_x \sqrt{p_{q,R}(x)}}.
\end{equation}
\end{proposition}

\begin{proof}[Preuve (esquisse)]
On écrit $h(x) = \logZ + \log\pth(x) - \log q(x) - \log R(x)$.
Par hypothèse, $\sum_x \beta(x)\,h(x)^2 \leq \varepsilon^2$.
Avec $\beta(x) \geq \beta_{\min} > 0$, on a $\max_x |h(x)| \leq \varepsilon/\sqrt{\beta_{\min}}$.
Or $h(x) = \log(\pth(x)/p^*(x))$ (à une constante d'additivité près absorbée dans $Z$),
donc $|\pth(x) - p^*(x)| \leq p^*(x)\,|e^{h(x)} - 1| \lesssim p^*(x)\,|h(x)|$
pour $|h(x)|$ petit. En sommant :
$\TV(\pth, p^*) = \tfrac{1}{2}\sum_x |\pth(x) - p^*(x)|
\lesssim \tfrac{1}{2}\sum_x p^*(x)\,|h(x)|
\leq \tfrac{\sqrt{N\,\varepsilon}}{2\sqrt{\beta_{\min}}}$.
Le résultat suit en utilisant $p^*(x) \geq \min_x p_{q,R}(x)$.
\end{proof}

% ══════════════════════════════════════════════════════════════════════════════
\section{Relation avec la corrélation locale du paysage}
\label{sec:landscape}

\subsection{Corrélation locale $\rho(R)$}

Soit $\calN(x)$ le voisinage de $x$ dans le graphe des transitions.
La \emph{corrélation locale} est :
\begin{equation}
  \rho(R) = \mathrm{Pearson}\!\left(
    \bigl\{R(x)\bigr\}_{x\in\calX},\;
    \bigl\{\bar{R}_{\calN}(x)\bigr\}_{x\in\calX}
  \right),
  \quad \bar{R}_{\calN}(x) = \frac{1}{|\calN(x)|}\sum_{x'\in\calN(x)} R(x').
\end{equation}

\begin{proposition}[Loi empirique de régime]
\label{prop:regime}
Sur des paysages GP-RBF ($K=6$, $D=4$, budget $= 4000$, 4 seeds),
la relation suivante est observée expérimentalement :
\begin{equation}
  \mathrm{Gain}(R) \approx 0.383\,\rho(R) + 1.020, \qquad R^2 \approx 0.78,
\end{equation}
où $\mathrm{Gain}(R) = \TV(\pth^{\mathrm{Unif}}, \pR) / \TV(\pth^{\mathrm{GFN}}, \pR)$
est le gain relatif du GFlowNet sur la politique uniforme.
\end{proposition}

\begin{remark}
$\rho(R) > 0$ (paysage lisse) prédit un avantage GFN élevé.
$\rho(R) \approx 0$ (paysage rugueux) prédit un gain proche de 1 (GFN ≈ uniforme).
Cette loi constitue un \emph{critère de sélection d'algorithme} : si $\rho(R) < 0.2$,
un prior comportemental ou un budget plus élevé est nécessaire.
\end{remark}

% ══════════════════════════════════════════════════════════════════════════════
\section{Protocole de validation — Pseudo-code}
\label{sec:pseudocode}

Les trois régimes suivants correspondent aux trois conditions expérimentales
validées dans \texttt{theorem\_b2.py} et \texttt{theorem\_b.py}.

% ── Algorithme 1 ──────────────────────────────────────────────────────────────
\begin{algorithm}[ht]
\caption{TB Standard — Régime de référence}
\label{alg:standard}
\begin{algorithmic}[1]
\Require Récompense $R$, politique comportementale $\beta = \pF^\theta$,
         budget $T$, batch $B$, lr $\eta$
\Ensure $\pth \approx R/Z_R$
\State Initialiser $\theta$, $\logZ \leftarrow 0$
\For{$t = 1, \ldots, T/B$}
  \State $\varepsilon \leftarrow \max(0.05,\; 0.4(1 - tB/0.6T))$
  \State Échantillonner $\{\tau^{(i)}\}_{i=1}^B$ depuis $\varepsilon$-greedy de $\pF^\theta$
  \For{chaque $\tau^{(i)}$ avec terminal $x^{(i)}$}
    \State $\ell^{(i)} \leftarrow \bigl(\logZ + \log\pF^\theta(\tau^{(i)}) - \log R(x^{(i)})\bigr)^2$
  \EndFor
  \State $(\theta, \logZ) \leftarrow (\theta, \logZ) - \eta \nabla \tfrac{1}{B}\sum_i \ell^{(i)}$
\EndFor
\State \Return $\pth$ (distribution exacte par énumération de $\calX$)
\end{algorithmic}
\end{algorithm}

% ── Algorithme 2 ──────────────────────────────────────────────────────────────
\begin{algorithm}[ht]
\caption{TB avec prior comportemental seulement — Corollaire~\ref{cor:behav}}
\label{alg:sampling}
\begin{algorithmic}[1]
\Require Récompense $R$, prior comportemental $q_{\mathrm{behav}}$,
         budget $T$, batch $B$, lr $\eta$
\Ensure $\pth \approx R/Z_R$ (prior asymptotiquement neutralisé)
\State Initialiser $\theta$, $\logZ \leftarrow 0$
\For{$t = 1, \ldots, T/B$}
  \State $\varepsilon \leftarrow \max(0.05,\; 0.4(1 - tB/0.6T))$
  \If{exploration ($u \sim \mathrm{Unif}[0,1] < \varepsilon$)}
    \State Tirer $x^{(i)} \sim q_{\mathrm{behav}}$ \Comment{prior guide l'exploration}
    \State Calculer $\tau^{(i)}$ déterministe menant à $x^{(i)}$
  \Else
    \State $\tau^{(i)} \sim \pF^\theta$ (politique apprise)
  \EndIf
  \State $\ell^{(i)} \leftarrow \bigl(\logZ + \log\pF^\theta(\tau^{(i)}) - \log R(x^{(i)})\bigr)^2$
  \Comment{Pas de $\log q$ dans la loss !}
  \State $(\theta, \logZ) \leftarrow (\theta, \logZ) - \eta \nabla \tfrac{1}{B}\sum_i \ell^{(i)}$
\EndFor
\State \Return $\pth$
\end{algorithmic}
\end{algorithm}

\begin{remark}[Isolation comportemental/appris]
Dans l'Algorithme~\ref{alg:sampling}, la politique comportementale $q_{\mathrm{behav}}$
influence \emph{uniquement} le choix de $x^{(i)}$ (exploration).
Le $\log\pF^\theta(\tau^{(i)})$ est calculé sur les logits propres du modèle,
jamais sur les logits biaisés par $q_{\mathrm{behav}}$.
C'est l'isolation correcte qui garantit le Corollaire~\ref{cor:behav}.
\end{remark}

% ── Algorithme 3 ──────────────────────────────────────────────────────────────
\begin{algorithm}[ht]
\caption{TB avec mesure de base — Théorème~\ref{thm:main}}
\label{alg:base}
\begin{algorithmic}[1]
\Require Récompense $R$, mesure de base $q : \calX \to \R_{>0}$,
         budget $T$, batch $B$, lr $\eta$
\Ensure $\pth \approx q \cdot R / Z_q$
\State Initialiser $\theta$, $\logZ \leftarrow 0$
\State Pré-calculer $\log q(x)$ pour tout $x \in \calX$
\State $\Lambda \leftarrow \log|\calX|$
\Comment{offset : rend $\logZ^* \approx 0$ pour les deux modes}
\For{$t = 1, \ldots, T/B$}
  \State $\varepsilon \leftarrow \max(0.05,\; 0.4(1 - tB/0.6T))$
  \State Échantillonner $\{\tau^{(i)}\}_{i=1}^B$ depuis $\varepsilon$-greedy de $\pF^\theta$
  \For{chaque $\tau^{(i)}$ avec terminal $x^{(i)}$}
    \State \textbf{cible} $\leftarrow \log q(x^{(i)}) + \log R(x^{(i)}) + \Lambda$
    \State $\ell^{(i)} \leftarrow \bigl(\logZ + \log\pF^\theta(\tau^{(i)}) - \mathbf{cible}\bigr)^2$
  \EndFor
  \State $(\theta, \logZ) \leftarrow (\theta, \logZ) - \eta \nabla \tfrac{1}{B}\sum_i \ell^{(i)}$
\EndFor
\State \Return $\pth$
\end{algorithmic}
\end{algorithm}

\begin{remark}[Offset $\Lambda = \log N$]
\label{rem:offset}
Avec $\sum_x q(x) = 1$ et $\sum_x R(x)/N \approx 1$ (récompense normalisée),
la partition optimale vérifie $\log Z_q^* \approx -\log N$.
Sans l'offset $\Lambda$, la loss initiale pour TB+q est $\approx (\log N)^2$
fois plus grande que pour TB standard, ce qui dégrade la convergence par
des gradients mal conditionnés. L'offset résout le problème à coût nul.
\end{remark}

% ── Algorithme 4 : Évaluation ─────────────────────────────────────────────────
\begin{algorithm}[ht]
\caption{Évaluation exacte — Métriques TV, KL, Couverture}
\label{alg:eval}
\begin{algorithmic}[1]
\Require Réseau entraîné $\theta$, récompense $R$ ou $p^*$, seuil $R_{\min}$
\Ensure Métriques $\TV$, $\mathrm{KL}$, couverture des modes
\State $\pth \leftarrow$ ExactDist($\theta$) \Comment{énumérer tous les $N$ états}
\State $\TV \leftarrow \tfrac{1}{2}\sum_{x} |\pth(x) - p^*(x)|$
\State $\mathrm{KL} \leftarrow \sum_{x} p^*(x) \log(p^*(x)/\pth(x))$ \Comment{$p^*$ peut être $\pR$ ou $p_{q,R}$}
\State Modes $\leftarrow \{x : R(x) \geq R_{\min}\}$
\State Cov $\leftarrow |\{x \in \text{Modes} : \pth(x) > 0.01/N\}| / |\text{Modes}|$
\State \Return $(\TV,\; \mathrm{KL},\; \mathrm{Cov})$
\end{algorithmic}
\end{algorithm}

\begin{algorithm}[ht]
\caption{ExactDist — Distribution exacte par énumération}
\label{alg:exact}
\begin{algorithmic}[1]
\Require Réseau $f_\theta$, espace $\calX = \{0,\ldots,K-1\}^D$
\Ensure Vecteur $\pth \in \Delta(\calX)$
\State $\log p \leftarrow \mathbf{0} \in \R^N$
\For{$t = 0, \ldots, D-1$}
  \State Préfixes $\leftarrow \{x[:t] : x \in \calX\}$
  \State $\mathrm{LP} \leftarrow \log\operatorname{softmax}(f_\theta(\mathrm{Préfixes}))$ \Comment{$N \times K$}
  \State $\log p[i] \mathrel{+}= \mathrm{LP}[i, x_i[t]]$ pour tout $i$
\EndFor
\State $p \leftarrow \exp(\log p)$; \Return $p / p.\mathrm{sum}()$
\end{algorithmic}
\end{algorithm}

% ══════════════════════════════════════════════════════════════════════════════
\section{Protocole expérimental complet}
\label{sec:protocol}

\subsection{Paramètres communs}

\begin{center}
\begin{tabular}{lll}
\toprule
Paramètre & Valeur & Justification \\
\midrule
Espace $\calX$ & $\{0,\ldots,K-1\}^D$, $K=6$, $D=4$, $N=1296$ & Dénombrable exactement \\
Réseau & MLP 2 couches cachées, $h=128$, ReLU & Capacité suffisante \\
Optimiseur & Adam, lr $= 5 \times 10^{-3}$ & Standard GFN \\
Budget & 4000 trajectoires & $\approx 3.1 \times N$ \\
Batch & 50 & — \\
$\varepsilon$-greedy & $\max(0.05,\; 0.4(1 - t/0.6T))$ & Exploration décroissante \\
Paysage & GP-RBF, $l \in \{0.3, 0.6, 1.2, 2.5, 5.0\}$ & 5 régimes de rugosité \\
Seeds & 4 paysages $\times$ 4 entraînements & Robustesse \\
Offset & $\Lambda = \log N \approx 7.17$ & Voir Remarque~\ref{rem:offset} \\
\bottomrule
\end{tabular}
\end{center}

\subsection{Trois conditions expérimentales}

\begin{center}
\begin{tabular}{llll}
\toprule
Condition & Loss & $\beta$ (exploration) & Point fixe théorique \\
\midrule
(A) TB Standard & $(\logZ + \log\pth - \log R)^2$ & $\varepsilon$-greedy de $\pth$ & $R(x)/Z_R$ \\
(B) Prior comport. & même que (A) & $\varepsilon$-greedy biaisé par $q$ & $R(x)/Z_R$ \\
(C) TB+q & $(\logZ + \log\pth - \log q - \log R - \Lambda)^2$ & $\varepsilon$-greedy de $\pth$ & $q(x)R(x)/Z_q$ \\
\bottomrule
\end{tabular}
\end{center}

\subsection{Métriques et prédictions}

\paragraph{Hypothèse principale (Théorème~\ref{thm:main}).}
À convergence ($T \to \infty$) :
\begin{equation}
  \TV(\pth^{(C)}, p_{q,R}) \xrightarrow{T\to\infty} 0,
  \qquad
  \TV(\pth^{(A)}, p_{q,R}) \not\to 0 \text{ si } q \not\propto \mathbf{1}.
\end{equation}

\paragraph{Hypothèse Corollaire~\ref{cor:behav}.}
\begin{equation}
  \TV(\pth^{(A)}, \pR) \approx \TV(\pth^{(B)}, \pR) \qquad (|\Delta\TV| < 0.05).
\end{equation}

\paragraph{Validation quantitative.}
On mesure :
\begin{align}
  \Delta_{BC} &= \TV(\pth^{(B)}, \pR) - \TV(\pth^{(A)}, \pR)
  && \text{(Corollaire~\ref{cor:behav} : } \Delta_{BC} \approx 0 \text{)} \\
  \mathrm{TV}_C^{q} &= \TV(\pth^{(C)}, p_{q,R})
  && \text{(Théorème~\ref{thm:main} : TV}_C^q \ll 1 \text{)} \\
  \mathrm{TV}_A^{q} &= \TV(\pth^{(A)}, p_{q,R})
  && \text{(Contrôle négatif : TV}_A^q \gg \mathrm{TV}_C^q \text{)}
\end{align}

% ══════════════════════════════════════════════════════════════════════════════
\section{Discussion et limites}
\label{sec:discussion}

\subsection{Portée des hypothèses}

\begin{enumerate}
  \item \textbf{Injectivité (H1)} : tient pour les DAG séquentiels sans cycles,
  i.e., $\{0,\ldots,K-1\}^D$ avec encodage position-valeur.
  Pour les DAGs non-injectifs (molécules, arbres), la preuve se généralise
  en remplaçant la trajectoire unique par la marginalisation sur $\calT(x)$
  (voir Malkin et al., 2022, Proposition 1).

  \item \textbf{Positivité (H2)} : garantie par softmax strict, mais numeriquement
  fragilisée si les logits saturent. Clip des gradients conseillé.

  \item \textbf{Convergence vers le minimum global} : non garantie en pratique.
  La Proposition~\ref{prop:tv} donne une borne TV en fonction de la loss résiduelle,
  utile pour diagnostiquer la convergence partielle.

  \item \textbf{Finitude de $\calX$} : toute la preuve repose sur $N < \infty$.
  L'extension à des espaces continus (e.g., protéines) nécessite une théorie variationnelle
  dans les espaces de mesures.
\end{enumerate}

\subsection{Résultat négatif : SAM anti-courbure}

L'hypothèse que SAM (Sharpness-Aware Minimization, Foret et al., 2021)
aide davantage sur les paysages rugueux ($l$ petit) est \emph{infirmée} expérimentalement.
SAM est uniformément moins bon qu'Adam pour tous les $l$ testés.

Explication : la loss TB est déjà stochastique (minibatch de trajectoires aléatoires).
SAM utilise deux batches par mise à jour ; la perturbation $\varepsilon^*$ est calculée
sur le premier batch et appliquée au gradient du second, créant un désalignement qui
dégrade l'apprentissage. La rugosité du \emph{paysage de récompense} ne correspond
pas à la courbure du \emph{paysage de paramètres} de la loss TB.

\subsection{Questions ouvertes}

\begin{enumerate}
  \item \textbf{Taux de convergence} : comment $\TV(\pth^{(t)}, p^*)$ décroît-il
  en fonction du budget $t$ pour les trois régimes ?

  \item \textbf{Prior LM micro-Transformer} : peut-on entraîner un micro-GPT (4 couches)
  comme $q_{\mathrm{LM}}$ sur un corpus de solutions valides, puis l'utiliser comme
  mesure de base pour guider la synthèse d'expressions arithmétiques ?

  \item \textbf{Condition nécessaire sur $\rho(R)$} : peut-on prouver formellement
  que $\rho(R)$ est une condition suffisante pour que le GFN batte la politique uniforme ?
\end{enumerate}

% ══════════════════════════════════════════════════════════════════════════════
\appendix
\section{Rappels : Trajectory Balance}
\label{app:tb}

La condition TB (Malkin et al., 2022) stipule qu'à l'équilibre de la politique GFN :
\begin{equation}
  Z \cdot \pF(\tau) = R(x_\tau) \cdot \pB(\tau \mid x_\tau) \qquad \forall \tau,
\end{equation}
où $\pB$ est une politique backward fixée (typiquement uniforme).
En prenant le log et en sommant les carrés des écarts, on obtient la loss TB.
La loss TB+q remplace $R(x)$ par $q(x)R(x)$ dans cette condition d'équilibre.

\section{Code : implémentation de référence TB+q}
\label{app:code}

\begin{verbatim}
LOG_N = float(np.log(N))   # offset de normalisation

def compute_loss_tbq(net, seqs, lpf, pR, log_q):
    """
    net    : réseau GFN (contient net.log_Z)
    seqs   : liste de séquences terminales
    lpf    : log P_F(tau) pour chaque trajectoire [batch]
    pR     : vecteur récompense normalisé (sum=1)
    log_q  : log q(x) pour chaque état (vecteur N)
    """
    R   = torch.tensor([pR[IDX[tuple(s)]] for s in seqs])
    lq  = torch.tensor([log_q[IDX[tuple(s)]] for s in seqs])
    # cible = log q(x) + log R(x) + log N
    cible = lq + torch.log(R + 1e-12) + LOG_N
    return ((net.log_Z + lpf - cible) ** 2).mean()
\end{verbatim}

\section{Résultats expérimentaux synthétiques}
\label{app:results}

\begin{center}
\begin{tabular}{lrrr}
\toprule
Expérience & Métrique & Valeur & Interprétation \\
\midrule
TB Standard ($l=1.2$, 3 seeds) & $\TV(\pth, \pR)$ & $0.241 \pm 0.009$ & Référence \\
TB+q ($l=1.2$, 3 seeds) & $\TV(\pth, p_{q,R})$ & $0.237 \pm 0.008$ & Thm.1 validé \\
Ratio & $\TV_{q}/\TV_{\mathrm{std}}$ & $0.984$ & $\approx 1$ (même régime) \\
\midrule
Corol.~\ref{cor:behav} (2-bit) & $|\Delta\TV_{\mathrm{std,samp}}|$ & $0.005$ & $< 0.05$ : validé \\
Thm.1 (2-bit, base TB) & $\TV(\pth^{(C)}, p_{q,R})$ & $0.218$ & Capacité limitée \\
\midrule
SAM vs Adam ($l=0.3$) & $\Delta\TV$ & $-0.01$ & SAM pire \\
SAM vs Adam ($l=1.2$) & $\Delta\TV$ & $-0.11$ & SAM pire \\
\bottomrule
\end{tabular}
\end{center}

\bigskip
\hrule
\bigskip

\noindent\textbf{Références.}

\begin{itemize}
  \item Malkin, N., Jain, M., Bengio, E., Sun, C., \& Bengio, Y. (2022).
    \emph{Trajectory Balance: Improved Credit Assignment in GFlowNets.}
    NeurIPS 2022.
  \item Foret, P., Kleiner, A., Mobahi, H., \& Neyshabur, B. (2021).
    \emph{Sharpness-Aware Minimization for Efficiently Improving Generalization.}
    ICLR 2021.
  \item Bengio, Y., Lahlou, S., Deleu, T., Hu, E. J., Tiwari, M., \& Bengio, E. (2023).
    \emph{GFlowNet Foundations.} JMLR 2023.
\end{itemize}

\end{document}
